\BeleskeID{02-s039}
% element: 02-s039:e01
% element: 02-s039:e02
% element: 02-s039:d01

Grupa reda \(k\) predstavlja \(2^k\) kombinacija u kojima \(k\) promenljivih uzima oba stanja, a ostale ostaju stalne. Za funkciju sa \(n\) ulaza takvoj grupi odgovara član sa \(n-k\) literala. Grupa od jednog polja ne uklanja nijednu promenljivu, dok grupa od četiri pravilno raspoređena polja uklanja dve.

Nije dovoljno izabrati proizvoljna četiri ili osam susednih polja. Grupa mora imati pravilnu strukturu pravougaonika, kvadrata, kvadra, kocke ili odgovarajućeg hiperkvadra odnosno hiperkocke. Svaka promenljiva koja se u grupi menja mora biti obuhvaćena svim potrebnim kombinacijama. Grupa oblika slova L, na primer, ne predstavlja jedan takav član. Pravilna grupa može biti razdvojena preko naspramnih ivica crteža.
